%% Master CV - Julian Askøe Bluming
%% Comprehensive overview of all competencies, experience, and achievements.
%% Use as a master reference when tailoring role-specific CVs.
%%
%% This master file may run past 2 pages; TAILORED CVs must be exactly 2 pages.
%% Compile with: lualatex main_example.tex
%% (pdflatex often fails on modern MiKTeX with fontawesome5 font-expansion errors.)

\documentclass[11pt,a4paper,sans]{moderncv}
\moderncvstyle{banking}
\moderncvcolor{blue}

% Force the name and section headings to render in moderncv blue (color1).
% Default banking leaves them black: moderncvstylebanking.sty's \colorlet
% copies (not aliases) the pre-scheme accent colour, so the name colours are
% frozen before \moderncvcolor runs. Re-let them after. \namefont is the hook
% every name-style macro routes through, so this also works on moderncv 2.3.1
% (Debian/Ubuntu apt), which has no \firstnamestyle/\lastnamestyle at all.
\renewcommand*{\namefont}{\fontsize{34}{36}\bfseries\upshape}
\colorlet{firstnamecolor}{color1}
\colorlet{lastnamecolor}{color1}
\colorlet{namecolor}{color1}
\renewcommand*{\sectionstyle}[1]{{\sectionfont\color{color1}#1}}

\usepackage[utf8]{inputenc}
% T1 is required for a Danish CV: it puts æ ø å in the font as single glyphs,
% so the PDF text layer extracts them as real characters for ATS parsers.
% Without it, OT1 composes them from accents and extraction degrades.
\usepackage[T1]{fontenc}
% Danish hyphenation via babel is OPTIONAL and NOT enabled: this machine has no
% babel-danish (Arch: `pacman -S texlive-langeuropean`; TeX Live: `tlmgr install
% hyphen-danish babel-danish`). Add \usepackage[danish]{babel} once installed --
% it improves line breaking only. The æ ø å glyphs come from fontenc T1 above and
% work without it.
% moderncv loads hyperref itself in an \AtEndPreamble hook, so \hypersetup
% must go in an \AtEndPreamble of our own: on moderncv < 2.4 a top-level
% \usepackage{hyperref} clashes with the class's own
% \RequirePackage[unicode]{hyperref}. From 2.4.0 the class passes its options
% through \PassOptionsToPackage instead, which is what removes that clash.
\AtEndPreamble{\hypersetup{
    colorlinks=true,
    linkcolor=blue,
    filecolor=magenta,
    urlcolor=blue,
    pdftitle={Julian Askoe Bluming - CV},
    % Keep pdfpagemode=UseNone: this block runs after moderncv's own
    % \AtEndPreamble (moderncv.cls sets pdfpagemode there), so a FullScreen
    % value here would win and open every CV in fullscreen presentation mode.
    pdfpagemode=UseNone,
}}
\usepackage[scale=0.80]{geometry}
\usepackage{needspace}
\usepackage{import}

% personal data
\name{Julian}{Askøe Bluming}
\address{Pansvej 12, 2680 Solrød Strand, Danmark}{}{}
\phone[mobile]{+45 91 56 85 28}
\email{julian.askoe@gmail.com}
% No LinkedIn URL on file. When one exists, replace the line below with:
%   \extrainfo{Autoriseret klinisk diætist \textbar{} \href{https://www.linkedin.com/in/...}{LinkedIn}}
\extrainfo{Autoriseret klinisk diætist}

\begin{document}

\makecvtitle

% ============================================================
%     PROFIL
% ============================================================

\vspace{6pt}
\small{Autoriseret klinisk diætist med kandidatgrad i Human Ernæring fra Københavns Universitet (forsvaret aug.\ 2026). Kombinerer klinisk diætetisk patientbehandling med praktisk erfaring i den formelle forskningsproces: fra protokoludarbejdelse og VEK-godkendelse til dataindsamling, databehandling og statistisk analyse. 20 ugers klinisk praktik fordelt på nefrologi, kardiologi, neurologi og endokrinologi. Udpræget detaljeorienteret med høj kvalitetssans i datadokumentation, selvkørende i egne opgaver og bevidst opsøgende på faglig sparring.}

% ============================================================
%     KERNEKOMPETENCER
% ============================================================

\section{Kernekompetencer}
\vspace{1pt}
\begin{itemize}

\item \textbf{Klinisk diætetik}: Diætetisk vejledning og ernæringsterapi inden for nefrologi (fosfat-, kalium- og proteinregulering, prædialyse og dialyse), kardiologi (hjertesund og kolesterolsænkende kost), endokrinologi (type 1/2-diabetes, pancreasresektion), neurologi og onkologi.

\item \textbf{Ernæringsvurdering}: Kostregistrering og -beregning, ernæringsscreening og risikovurdering, dækningsgradsberegning, Nutrition Impact Symptoms (NIS-faktorer), perioperativ og ERAS-ernæring.

\item \textbf{Klinisk forskningsmetode}: Protokol- og deltagerinformationsudarbejdelse, VEK-ansøgning, registrering på ClinicalTrials.gov, samtykkeindhentning, deltagerrekruttering og -opfølgning, systematisk dataindsamling og kvalitetssikring.

\item \textbf{Data og analyse}: REDCap, VitaKost, GraphPad Prism, R-studio (kursusniveau), journal- og lab-dataudtræk fra Sundhedsplatformen.

\item \textbf{Evidensvurdering}: Systematisk litteratursøgning (PICO), kritisk analyse og vurdering af studiedesign og forskningslitteratur.

\item \textbf{Formidling og samarbejde}: Individuel patientvejledning og undervisning, udarbejdelse af diætetisk vejledningsmateriale, tværfagligt samarbejde med læger og sygeplejersker på afdelingsniveau.

\end{itemize}

% ============================================================
%     FORSKNINGSERFARING
% ============================================================

\section{Forskningserfaring}
\vspace{3pt}
\begin{itemize}

\needspace{5\baselineskip}
\item{\cventry{2024-2026}{Specialestuderende, forskningsansvarlig}{Københavns Universitet}{København}{}{\vspace{1pt}
\begin{itemize}
    \item {Klinisk tillægsstudie (n=42) til en større kohorte: korrelation mellem nyrefunktion (clearance) og calcium-/fosfatbalance hos stadie 4-5 kronisk nyresyge patienter, samt metodesammenligning af kostindtag via vejet registrering vs.\ billedbaseret fotometode.}
    \item {Udarbejdede udkast til tillægsprotokol og deltagerinformation, godkendt af Videnskabsetisk Komité (VEK), og registrerede studiet på ClinicalTrials.gov.}
    \item {Varetog participant-information, samtykkeindhentning og oplæring af deltagere i kostregistrering i koordination med forsøgslæge, samt løbende og telefonisk diætetisk opfølgning.}
    \item {Udtræk af lab- og journaldata fra Sundhedsplatformen, systematisk databearbejdning i VitaKost og REDCap, statistisk analyse i GraphPad Prism.}
    \item {Udarbejdede diætetisk vejledningsmateriale til studiet. Forbereder aktuelt studiet til videnskabelig publicering som førsteforfatter.}
\end{itemize}}}

\vspace{3pt}

\needspace{5\baselineskip}
% Exact months of the bachelor project are not on file; a lone year imports
% worse into ATS forms than a range. Replace with e.g. {feb. 2024 - jun. 2024}
% once confirmed - do not guess.
\item{\cventry{2024}{Bachelorprojekt}{Rigshospitalet, Ernæringsenheden}{København}{}{\vspace{1pt}
\begin{itemize}
    \item {Bidrog til opstartsfasen og den systematiske dataindsamling i et prospektivt, observationelt studie af ERAS-behandlede cystektomi-patienter på urologisk afdeling (n=25).}
    \item {Kortlægning af energi- og proteinindtag samt dækningsgrad (>75\%), registrering af Nutrition Impact Symptoms, opfølgning på vægtændringer og væskebalance over et 30-dages postoperativt forløb.}
\end{itemize}}}

\end{itemize}

% ============================================================
%     KLINISK PRAKTIK
% ============================================================

\section{Klinisk praktik}
\vspace{3pt}
\begin{itemize}

\needspace{5\baselineskip}
\item{\cventry{nov. 2023 - jan. 2024}{Klinisk praktik, nefrologisk afdeling (10 uger)}{Rigshospitalet}{København}{}{\vspace{1pt}
\begin{itemize}
    \item {Klinisk vejledning af prædialyse- og dialysepatienter inden for fosfat-, kalium- og proteinregulering.}
    \item {Vejledning og håndtering af endokrinologiske tilstande (type 1/2-diabetes og pancreasresekerede patienter).}
    \item {Løbende tværfaglig dialog og sparring med afdelingens læger og sygeplejersker.}
\end{itemize}}}

\vspace{3pt}

\needspace{5\baselineskip}
\item{\cventry{sep. 2023 - nov. 2023}{Klinisk praktik, hjerteafdelingen (6 uger)}{Bispebjerg og Frederiksberg Hospital}{København}{}{\vspace{1pt}
\begin{itemize}
    \item {Vejledning i hjertesund og kolesterolsænkende kost til patienter med hyperkolesterolæmi samt højproteinkost til postoperative patienter.}
    \item {Selvstændig undervisning i vægttab for knæartrosepatienter i forbindelse med et forskningsprojekt i den reumatologiske forskningsenhed (1 ugedag).}
\end{itemize}}}

\vspace{3pt}

\needspace{5\baselineskip}
\item{\cventry{aug. 2023 - sep. 2023}{Klinisk praktik, neurologisk afdeling (4 uger)}{Rigshospitalet Glostrup}{Glostrup}{}{\vspace{1pt}
\begin{itemize}
    \item {Deltog i opstart af ernæringsterapi og diætetisk praksis over for neurologiske patienter (bl.a.\ tetraplegi og Parkinson), spiseforstyrrelser og diabetes.}
\end{itemize}}}

\end{itemize}

% ============================================================
%     UDDANNELSE
% ============================================================

\section{Uddannelse}
\vspace{1pt}
\begin{itemize}

\needspace{5\baselineskip}
\item{\cventry{2024-2026}{Kandidat (MSc) i Human Ernæring}{Københavns Universitet}{København}{}{\vspace{1pt}
Kandidatspeciale: klinisk tillægsstudie af nyrefunktion og calcium-/fosfatbalance ved stadie 4-5 kronisk nyresygdom samt metodesammenligning af kostregistrering (se Forskningserfaring).\\
Relevante kurser: \emph{Study Design in Human Nutrition} (protokoludarbejdelse, studiedesign, statistisk analyse i R-studio); \emph{Evidence, Diet and Health} (kritisk analyse af forskningslitteratur og studiedesigns); \emph{Cell Cycle and Cancer Control} (molekylær onkologi).
}}

\vspace{3pt}

\needspace{5\baselineskip}
\item{\cventry{2021-2024}{Professionsbachelor i Ernæring og Sundhed (Klinisk Diætetik)}{Københavns Professionshøjskole}{København}{}{\vspace{1pt}
Bachelorprojekt gennemført i Ernæringsenheden på Rigshospitalet (se Forskningserfaring).\\
Relevant kursus: \emph{Tolkning og anvendelse af klinisk forskning} (systematisk litteratursøgning efter PICO, kritisk vurdering af kliniske studier).
}}

\end{itemize}

% ============================================================
%     AUTORISATION OG KURSER
% ============================================================

\section{Autorisation}
\vspace{1pt}
\begin{itemize}
\item {Autoriseret klinisk diætist (dansk autorisation).}
\item {Kørekort under erhvervelse, forventet erhvervet nov.\ 2026.}
\end{itemize}

% ============================================================
%     SPROG
% ============================================================

\section{Sprog}
\vspace{1pt}
\begin{itemize}
\item {Dansk (modersmål), engelsk (godt arbejdsniveau i skrift og faglitteratur).}
\end{itemize}

% ============================================================
%     PUBLIKATIONER
% ============================================================

\section{Publikationer}
\vspace{1pt}
\begin{itemize}
\item {Bluming, J.\ et al.\ Klinisk tillægsstudie af nyrefunktion og calcium-/fosfatbalance hos stadie 4-5 kronisk nyresyge patienter samt metodesammenligning af kostregistrering. \emph{Manuskript under udarbejdelse til publicering, førsteforfatter.}}
\end{itemize}

% ============================================================
%     INTERESSER
% ============================================================

\section{Interesser}
\vspace{1pt}
\begin{itemize}
\item {Ernæring og styrketræning.}
\end{itemize}

% ============================================================
%     REFERENCER
% ============================================================

\section{Referencer}
\vspace{1pt}
\begin{itemize}
\item {Jens Rikardt Andersen, MD, Københavns Universitet -- specialevejleder. jra@nexs.ku.dk, +45 23 34 66 54.}
\item {Kirsten Thal Jantzen, klinisk diætist, Herlev Hospital -- vejleder, nefrologisk afdeling. kirsten.thal-jantzen@regionh.dk, +45 60 75 76 18.}
\item {Janne Christensen, klinisk diætist, Rigshospitalet Glostrup -- praktikvejleder. janne.christensen@regionh.dk, +45 38 63 22 27.}
\item {Yderligere referencer udleveres på forespørgsel.}
\end{itemize}

\end{document}
